% Corrected French for printed p.199.
\clearpage
\SourcePage{199}

\LayerTarget{tohoku:construction:canonical-map-inverse-direct:p1}
$g(x)$ sur $V$, dont la valeur en chaque $x\in V$ est un élément de
\[
f_*^G(A)(f(x))=
\varinjlim_{U\ni f(x)}\Gamma(f^{-1}(U),A)^G,
\]
et telle que pour tout $x\in V$ existe un voisinage ouvert $V'\subset V$ de
$x$ et une section $h$ de $f_*^G(A)$ sur un voisinage $U'$ de $f(x)$ (i.e.
une section invariante $h$ de $A$ sur $f^{-1}(U')$) telle qu'on ait
$g(x)=h(x)$ pour $x\in V'\cap f^{-1}(U')$. On en conclut un monomorphisme
canonique\textsuperscript{10)} :

\LayerTarget{tohoku:equation:5.1.2}
\[
f^{-1}(f_*^G(A))\longrightarrow A.
\tag{5.1.2}
\]

\LayerTarget{tohoku:criterion:inverse-direct-essential-image}
identifiant $f^{-1}(f_*^G(A))$ à un sous-faisceau $A'$ de $A$. Il résulte de
la formule (5.1.1) ci-dessus que l'on a $A'\simeq A$ si et seulement si $A$ est
isomorphe à un $G$-faisceau de la forme $f^{-1}(B)$. S'il en est ainsi alors
pour tout $x\in X$, les opérations de $G_x$ dans $A(x)$ sont triviales. La
réciproque est d'ailleurs vraie si $G$ satisfait la condition (D) plus bas
(cf. 5.3), comme on le vérifie facilement ;

\LayerTarget{tohoku:result:fixed-point-free-equivalence}
\textit{si donc en plus les $G_x$ sont réduits à l'identité} (i.e. si aucun
$g\ne e$ dans $G$ n'admet de point fixe) \textit{les foncteurs $f_*^G$ et
$f^{-1}$ établissent des isomorphismes réciproques l'un de l'autre de la
catégorie des $G$-faisceaux sur $X$ et de la catégorie des faisceaux sur $Y$.}

\LayerTarget{tohoku:example:constant-direct-image}
\noindent\textsc{Exemples d'images directes de $G$-faisceaux.} Si $A$ est le
faisceau constant sur $X$ défini par un ensemble $M$ où $G$ opère
trivialement, alors $A^G$ est le faisceau constant sur $Y$ défini par le même
ensemble $M$ ; mais si $G$ n'opère pas trivialement dans $M$, et si on n'est
pas dans le cas ``sans points fixes'', alors $f_*^G(M)$ n'est évidemment plus
localement constant en général.

\LayerTarget{tohoku:example:germ-function-direct-images}
Supposons que la condition (D) de 5.3 soit vérifiée, soit $A$ le faisceau des
germes d'applications de $X$ dans un ensemble $M$, alors $A^G$ est le faisceau
des germes d'applications de $Y$ dans $M$. Dans cette correspondance, si $X$
est une variété différentiable (resp. un espace analytique réel, resp. un
espace analytique complexe) et si $A$ est le faisceau des germes de fonctions
(à valeurs dans un espace de même nature) différentiables (resp. analytiques
réelles, resp. analytiques complexes) alors $A^G$ est le faisceau de même nom
relatif à $Y$. Ce dernier exemple est particulièrement important, et l'énoncé
analogue est vrai (par définition d'ailleurs, comme pour les précédents!) en
Géométrie Algébrique abstraite, ou même pour les variétés arithmétiques.

\LayerTarget{tohoku:section:5.2}
\bigskip
\noindent\textbf{5.2. Les foncteurs $H^n(X;G,A)$ et $H^n(G,A)$ et les suites
spectrales fondamentales.}

\LayerTarget{tohoku:definition:section-invariant-functors}
Pour éviter des confusions, nous notons $\Gamma_X$ et $\Gamma_Y$ les foncteurs
``sections'' définis pour des faisceaux sur $X$ resp. sur $Y$, et $\Gamma^G$
le foncteur $M\longmapsto M^G$ faisant correspondre, à un ensemble $M$ où
opère $G$, l'ensemble des invariants de $M$. Enfin, si $A$ est un $G$-faisceau
sur $X$, on pose

\LayerTarget{tohoku:equation:5.2.1}
\[
\Gamma_X^G(A)=\Gamma_X(A)^G,
\tag{5.2.1}
\]
donc par définition, on a les formules

\LayerTarget{tohoku:equation:5.2.2}
\[
\Gamma_X^G=\Gamma^G\Gamma_X=\Gamma_Yf_*^G.
\tag{5.2.2}
\]

\LayerTarget{tohoku:footnote:10}
\begin{quote}
\footnotesize
10) Plus généralement, si $f$ est une application continue d'un espace $X$
dans un espace $Y$, on a, pour les faisceaux $A$ sur $X$, des monomorphismes
fonctoriels $f^{-1}(f_*(A))\longrightarrow A$, dont le monomorphisme (5.1.2)
se déduit immédiatement, grâce à l'injection $A^G\longrightarrow f_*(A)$.
\end{quote}

\TohokuContinuation{200}
